\clearpage
\fancyhead[L]{\small CONNAISSANCE DES TEMS.}
\SourceLeaf{614}
\begin{center}\small Göttingische gelehrte Anzeigen. Stück 40. 41. Seite 402..408. 1819 März 11.\\[3pt]\rule{.72\linewidth}{.3pt}\end{center}

\noindent\emph{Connaissance des tems ou des mouvemens célestes à l'usage des astronomes et des navigateurs pour l'an {\upshape 1820}, publiée par le bureau des longitudes.} Paris 1818. Bei Witwe Courcier. 446 Seiten. gr.~8.

In der Einrichtung des Kalenders finden wir bei diesem Jahrgange nur die Eine Abänderung, dass die Heiligen-Namen, welche sonderbar genug, seit ein paar Jahren zwischen die astronomischen Rubriken gesetzt waren, wieder weggeblieben sind; den dadurch gewonnenen Platz hat man dazu verwandt, beim Abstande des Widderpunktes von der Sonne und bei der Declination der Sonne eine Differenzcolumne beizufügen, statt deren wohl etwas Brauchbareres hätte gewählt werden können. Bei der Tafel für die geographische Lage der vornehmsten Punkte der Erde wird versichert, dass sie von neuem nachgesehen sei, und dass man diese neueste Ausgabe als die genaueste und zuverlässigste anzusehen habe. Allein diese Versicherung ist ein stehender Artikel, und es bedarf bei diesem Bande wie bei seinen Vorgängern, eben keines langen Suchens, um sich zu überzeugen, dass dies Verzeichniss mit vieler Nachlässigkeit redigirt wird. So finden wir z.\,B. die Polhöhen von Ofen um $32''$. die von Göttingen um $24' 26''$ zu gross, die von Königsberg um $38''$ zu klein, die von Mannheim um $5''$ zu gross: bei dem ersten und dritten dieser Oerter mögen andere Punkte als die gegenwärtigen Sternwarten zu Grunde liegen; dies hätte aber wenigstens nicht sein sollen. Von den Zusätzen berühren wir hier nur die Originalaufsätze, indem wir die auch bei diesem Bande sehr zahlreichen weitläuftigen Auszüge aus andern Büchern und Abhandlungen mit Stillschweigen übergehen. Den Anfang machen die schätzbaren Tafeln der Vesta von \textsc{Daussy}, die sich auf die im Jahrgange 1818 abgedruckten Elemente und Störungsformeln gründen. Hierauf folgt eine lehrreiche Abhandlung über den Apparat zur Bestimmung der Länge des Secundenpendels, von \textsc{Laplace}, worin besonders der Einfluss des Umstandes untersucht wird, dass die Schneide, auf der das Pendel schwingt, nicht in aller Schärfe eine mathematische Linie, sondern ein Stück einer kleinen Cylinderfläche ist. Die vortrefflichen \textsc{Borda}'schen Versuche über diesen Gegenstand verdienten diese subtile Berücksichtigung. --- Ueber die Bestimmung des Perpendikels auf den Erdmeridian, und über verschiedene andre sich darauf beziehende Aufgaben von \textsc{Puissant}. Nach unserm Dafürhalten ist die ganze bisherige Behandlungsart dieser Gegenstände noch nicht die rechte; und erst, wenn man dieselben aus neuen Gesichtspunkten betrachten wird (was hier nicht geschehen ist), darf man hoffen, Vollständigkeit und Schärfe mit Einfachheit und Geschmeidigkeit in den Rechnungsoperationen zu vereinigen. Es scheint um so wünschenswerther, dass die Geometer hierauf ihre Aufmerksamkeit richten, da gegenwärtig von mehrern Regierungen grosse Messungen veranstaltet werden, die, mit den besten Hülfsmitteln ausgeführt, neue und wichtige Aufschlüsse über die Gestalt der Erde hoffen lassen. Nach einer von Hrn. \textsc{Puissant} hier gegebenen Nachricht hat auch das Französische Gouvernement eine neue Triangulirung beschlossen, die an die Mittagslinie von Dünkirchen und an das Perpendikel von Brest nach Strassburg geknüpft und über das ganze Königreich ausgedehnt werden soll. --- Ueber die elliptische Bahn des Cometen von 1783 und ihre Aehnlichkeit mit der Bahn des Cometen von 1793, von \textsc{Burckhardt}. Die geringe Uebereinstimmung der Beobachtungen des Cometen von 1783 mit den von verschiedenen Astronomen versuchten parabolischen Elementen veranlasste Hrn. \textsc{Burckhardt}, diese Bahn als elliptisch zu berechnen. Hr. \textsc{Burckhardt} brachte eine Ellipse mit einer Umlaufszeit von $5\tfrac{1}{2}$ Jahren heraus, welche die Beobachtungen bei weitem besser darstellt. Auch eine Ellipse von 10 Jahren gab noch eine ganz gute Uebereinstimmung. Da diese Elemente einige Aehnlichkeit mit denen des zweiten Cometen von 1793 haben, so schien es denkbar, dass beide Cometen identisch wären, wesshalb \SourceLeaf{615}Hr. \textsc{Burckhardt} auch die Beobachtungen des Cometen von 1793, welche bisher noch nicht gedruckt waren, einer neuen Berechnung unterwarf. Leider sind diese Beobachtungen wenig genau. Es ist nicht wol möglich, nach den hier gelieferten Resultaten ein bestimmtes Urtheil zu fällen, zumal da wahrscheinlich Druckfehler dabei eingeschlichen sind, deren Verbesserung sich nicht errathen lässt (die Länge des Knoten, welche \textsc{Saron} in einer parabolischen Bahn $83^\circ\,55'$ gefunden hatte, wird von Hrn. \textsc{Burckhardt}, in der einen elliptischen Bahn gesetzt $359'\,4'\,48''$ (sic), und bei einer zweiten elliptischen Bahn ganz ausgelassen). Inzwischen verdient Hr. \textsc{Burckhardt} Dank, die Aufmerksamkeit der Astronomen auf diesen Cometen gelenkt und die \textsc{Messier}'schen Beobachtungen desselben hier zuerst bekannt gemacht zu haben. --- Auszug aus einem Schreiben des Hrn. \textsc{Schumacher} an Hrn. \textsc{Delambre}. Die Unterschiede zwischen \textsc{Tycho}'s und \textsc{Picard}'s Azimuthalbestimmungen auf der Insel Hueen werden hier aus einer Abhandlung \textsc{Augustin}'s befriedigend aufgeklärt. Von Helsingborg hatte \textsc{Picard} nicht den rechten Thurm, sondern einen erst nach \textsc{Tycho}'s Zeit erbauten beobachtet, und übrigens hatte \textsc{Tycho} anfangs, mit unvollkommnen Instrumenten, seine Mittagslinien um einen Viertelsgrad falsch gezogen, welchen Fehler er aber schon selbst einige Jahre später zufolge der Historia Coelestis bemerkt und berichtigt hatte. --- Ueber einen Aufsatz des Hrn. \textsc{Littrow} über die thermometrische Correction der Strahlenbrechung. Hr. \textsc{Delambre} rechtfertigt hier sehr weitläuftig seine eigne Tafel, die Hr. \textsc{Littrow} der Ungenauigkeit beschuldigt hatte. Allein obgleich diese Rechtfertigung 10 Seiten einnimmt, findet man doch nicht klar hervorgehoben, worin eigentlich Hrn. \textsc{Littrow}'s Versehen bestand. Es war dieses, dass Hr. \textsc{Littrow} durch Uebereilung vorausgesetzt hatte, wenn die Dichtigkeit der Luft, die beim Gefrierpunkt als Einheit angenommen für die Temperatur von $\alpha$ Grad über denselben durch die Formel 
\[
\frac{1}{1+0.0046875\,\alpha}
\]
 ausgedrückt wird, \emph{dieselbe} Formel gültig bleibe, wenn die Dichtigkeit für $8^\circ$ als Einheit betrachtet wird, und dann $\alpha$ den Thermometerstand über $8^\circ$ bedeutet, welches offenbar unrichtig ist. --- Ueber ein neues Mittel, die Dauer der Pendelschwingungen zu reguliren, von Hrn. \textsc{Prony}. Die Pendelstange ist etwas nach oben verlängert; an derselben ist eine sehr dünne kurze Querstange auf Reibung so befestigt, dass sie sich in einer auf jener senkrechten Ebene drehen lässt; an jedem Ende trägt diese Querstange ein Kügelchen von Platina. Offenbar wird das Moment der Trägheit dieses Pendels, in Beziehung auf die Schwingungsaxe, am kleinsten, wenn die Querstange mit dieser Axe parallel ist, am grössten hingegen, wenn beide einen rechten Winkel mit einander machen. Im ersten Fall werden daher die Schwingungen am schnellsten, im andern am langsamsten. Bei schicklich gewählten Dimensionen des Apparats ist hiedurch eine sehr feine Regulirung möglich, wobei ein nicht unwichtiger Vortheil der ist, dass man den Gang corrigiren kann, ohne die Uhr anzuhalten. --- Ueber mehrere unter den \textsc{Flamstead}'schen Sternen gefundene Beobachtungen des Planeten Uranus, von Hrn. \textsc{Burckhardt}. Ganz zufällig fand Hr. \textsc{Burckhardt} bei Gelegenheit anderer Nachsuchungen in \textsc{Flamstead}'s Historia Coelestis noch \emph{fünf} Uranusbeobachtungen (1712 März 22; 1715 Febr. 21. 22. 25 und April 18, alles nach altem Styl). Dieser Fund ist um so unerwarteter, da alle Astronomen diesen Gegenstand durch Hrn. \textsc{Bode}'s Nachforschung erschöpft glaubten, und um so schätzbarer, weil die Beobachtungen von 1715 eine sehr gute Bestimmung der Opposition liefern, während die isolirte Beobachtung von 1690 noch einer Ungewissheit ausgesetzt sein kann. \textsc{Flamstead} hat also unbewusster Weise den Planeten nicht weniger als sechsmal beobachtet, und die Entdeckung hätte ihm selbst gar nicht entgehen können, wenn er die Gewohnheit gehabt hätte, seine Beobachtungen immer gleich auf der Stelle wenn auch nur oberflächlich zu revidiren, und unter sich zu vergleichen. --- Ueber die kleinen Gleichungen in der Jupitersbewegung von \emph{ebendemselben}. Hr. \textsc{Burckhardt} gibt die numerischen Werthe mehrerer von \textsc{Laplace} übergangener Gleichungen der \SourceLeaf{616}zweiten bis sechsten Ordnung, unter denen einige doch nicht ganz unerheblich sind. Da die ähnlichen in der Theorie der Saturnsbewegung übergangenen Gleichungen noch beträchtlich grösser sein müssen, so wäre es wohl möglich, dass dieser Umstand mit zu dem auffallenden Unterschiede beigetragen hat, der bekanntlich bei Bestimmung der Jupitersmasse aus der Saturns- und aus der Pallasbewegung sich ergeben hat. --- Von \emph{demselben}, Bemerkungen über verschiedene Fixsterne. --- Cometenbeobachtungen auf der Königl. Sternwarte in Paris von 1811---1815 (grösstentheils schon anderwärts gedruckt), nebst Elementen, die durch Hrn. \textsc{Nicollet} daraus abgeleitet sind. --- Anwendung der Wahrscheinlichkeitsrechnung auf die geodätischen Operationen von \textsc{Laplace}. Die Wahrscheinlichkeitsrechnung findet da ihre Anwendung, wo mehr als das unumgänglich nöthige beobachtet ist, wenn die Beobachtungen keine absolute Genauigkeit haben. Sie lehrt die vortheilhafteste Combination der Beobachtungen und die Bestimmung der relativen Zuverlässigkeit der Resultate. Hr. \textsc{Laplace} betrachtet hier den Fall, wo auf der Oberfläche einer Kugel zwei Punkte durch ein Netz von Dreiecken verbunden, in den Dreiecken alle Winkel beobachtet, und \emph{zwei} Grundlinien, eine zu Anfang, die andere beim Ende des Netzes gemessen sind. Die Winkelmessungen werden als mit Fehlern behaftet, die Grundlinie als fehlerfrei betrachtet, und gesucht werden die Länge des die zwei Punkte verbindenden Bogens, und die Winkel, welche er mit der ersten und letzten Dreiecksseite macht. Man sieht, dass diese Aufgabe nur sehr speciell, und die Anwendung der Wahrscheinlichkeitstheorie auf höhere Geodäsie dadurch keinesweges erschöpft ist.
